\section*{Limitations}
\label{sec:limitations}

Several limitations qualify the system. First, \textbf{live-server
flakiness}: LIVE mode depends on third-party MCP servers that can be slow or
unreachable, which is precisely why REPLAY is the default and every LIVE
stage falls back to a frozen fixture; the trade-off is that the most
reliable path through the demo is not, strictly, live. Second,
\textbf{scorer conservatism}: the studio inherits the methodology's
deliberately conservative effect scorer, which accepts a run only when the
required effects are demonstrably present and can therefore fail an otherwise
valid alternative trajectory; the verdicts shown are lower bounds on
capability, not a definitive judgment of an agent. Third, \textbf{a
single-task default showcase}: the booth experience centers on one worked
example and three curated fixtures chosen to make the verdict flip vivid, so
it illustrates the disagreement rather than measuring its prevalence (the
parent study does the latter). Fourth, \textbf{Tier-2 non-determinism}: the
optional effect-equivalence judge is an LLM and is therefore kept off the
demo's critical path; the headline verdict is Tier-1 only. Finally,
\textbf{the demo is a front-end, not new science}: it visualizes and wraps
the published pipeline and introduces no new scoring method; its
contribution is interactional and infrastructural, not methodological.
